\chapter{Phương pháp đề xuất}
\label{ch:method}

Chương này trình bày chi tiết hệ thống hỏi đáp trên đồ thị tri thức thời gian tiếng Việt mà khóa
luận xây dựng. Trước hết, Mục~\ref{sec:overview} phác họa kiến trúc tổng thể hai pha kế thừa từ
khung ARI. Mục~\ref{sec:cronqvn} mô tả quy trình xây dựng bộ dữ liệu CronQ-VN -- đóng góp thứ nhất.
Mục~\ref{sec:knowledge-based} và Mục~\ref{sec:knowledge-agnostic} lần lượt trình bày pha tương tác
dựa trên tri thức và pha suy luận trừu tượng. Mục~\ref{sec:vietnamese-adapt} trình bày các thích
nghi cho tiếng Việt và vận hành cục bộ -- đóng góp thứ hai. Cuối cùng, Mục~\ref{sec:crossencoder}
đề xuất bộ lọc hành động dựa trên cross-encoder -- đóng góp thứ ba.

\section{Tổng quan kiến trúc hệ thống}
\label{sec:overview}

Hệ thống kế thừa khung Suy luận Trừu tượng Quy nạp (ARI)~\cite{chen2024ari}, chia quá trình suy
luận thời gian thành hai pha tách bạch nhằm giải quyết đồng thời hai hạn chế của mô hình ngôn ngữ
lớn đã nêu ở Mục~\ref{sec:limitations}.

\begin{itemize}
  \item \textbf{Pha tương tác dựa trên tri thức} (\textit{knowledge-based}) đảm nhận toàn bộ thao
  tác trên đồ thị tri thức: trích xuất đồ thị con, liệt kê các hành động nguyên tử khả thi, thực
  thi hành động và cập nhật trạng thái. Pha này cô lập mô hình ngôn ngữ khỏi khối tri thức cụ thể
  đồ sộ và nhiễu.
  \item \textbf{Pha suy luận trừu tượng} (\textit{knowledge-agnostic}) là nơi mô hình ngôn ngữ lớn
  đưa ra quyết định chiến lược ở mức cao: chọn hành động phù hợp dựa trên \emph{phương pháp luận
  trừu tượng} đã chắt lọc từ kinh nghiệm lịch sử, thay vì thao tác trực tiếp với tri thức cụ thể.
\end{itemize}

Toàn hệ thống vận hành theo hai quá trình (Hình~\ref{fig:architecture}). \textbf{Quá trình học}
(nét liền) duyệt một tập câu hỏi mẫu, ghi lại các trace suy luận, phân cụm chúng theo cấu trúc, rồi
chắt lọc một phương pháp luận trừu tượng cho mỗi cụm và lưu vào bộ nhớ. \textbf{Quá trình suy luận}
(nét đứt) với mỗi câu hỏi mới sẽ tìm cụm phù hợp nhất, lấy phương pháp luận tương ứng để dẫn dắt mô
hình ngôn ngữ lớn tương tác nhiều bước với pha tri thức cho tới khi đưa ra đáp án.

\begin{figure}[ht]
  \centering
  \fbox{\parbox{0.93\textwidth}{\centering\vspace{2mm}\footnotesize
  \textbf{[Pha tri thức]} Câu hỏi $\rightarrow$ liên kết thực thể $\rightarrow$ đồ thị con
  $\rightarrow$ liệt kê hành động $\rightarrow$ lọc $\rightarrow$ thực thi $\rightarrow$ cập nhật\\[2mm]
  $\updownarrow$ \emph{(mỗi bước)}\\[2mm]
  \textbf{[Pha trừu tượng]} Bộ nhớ phương pháp luận $\xrightarrow{\text{chọn cụm}}$ hướng dẫn
  $\rightarrow$ LLM chọn hành động\\[2mm]
  \emph{Học (nét liền): trace lịch sử $\rightarrow$ phân cụm $\rightarrow$ chắt lọc phương pháp luận}
  \vspace{2mm}}}
  \caption{Kiến trúc tổng thể của hệ thống: pha tương tác dựa trên tri thức (thực thi trên đồ thị)
  và pha suy luận trừu tượng (quyết định bằng mô hình ngôn ngữ lớn dưới hướng dẫn của phương pháp
  luận). \todo{vẽ lại bằng TikZ theo Figure~3 của~\cite{chen2024ari}, thể hiện rõ luồng học (nét
  liền) và luồng suy luận (nét đứt)}}
  \label{fig:architecture}
\end{figure}

\section{Xây dựng bộ dữ liệu CronQ-VN}
\label{sec:cronqvn}

Do chưa tồn tại bộ dữ liệu TKGQA tiếng Việt, khóa luận tự xây dựng CronQ-VN qua hai giai đoạn: (i)
trích xuất một đồ thị tri thức thời gian tiếng Việt từ Wikidata; (ii) sinh tự động các câu hỏi kèm
đáp án chuẩn từ đồ thị này. Quy trình lấy cảm hứng từ CronQuestions~\cite{saxena2021cronkgqa} và kỹ
thuật lọc của bài toán hoàn thiện đồ thị tri thức thời gian~\cite{lacroix2020tkbc}.

\subsection{Trích xuất đồ thị tri thức thời gian từ Wikidata}
\label{sec:buildkg}

Mô-đun \texttt{build\_kg.py} duyệt toàn bộ bản kết xuất (\textit{dump}) mới nhất của Wikidata theo
hai lượt:

\begin{enumerate}
  \item \textbf{Lượt 1 -- thu thập nhãn.} Duyệt mọi thực thể, thu thập ánh xạ từ mã định danh
  \texttt{QID} sang nhãn tiếng Việt (và nhãn tiếng Anh để tham chiếu) cho những thực thể có nhãn
  tiếng Việt.
  \item \textbf{Lượt 2 -- trích xuất sự kiện thời gian.} Với mỗi khẳng định, chỉ giữ lại sự kiện
  thỏa đồng thời: (a) chủ thể \emph{và} đối tượng đều có nhãn tiếng Việt; (b) khẳng định có ít nhất
  một định tố thời gian trong \{\texttt{P580}, \texttt{P582}, \texttt{P585}\}; và (c) mốc thời gian
  được rời rạc hóa về cấp năm.
\end{enumerate}

Sau hai lượt, đồ thị được lọc thưa theo tinh thần của~\cite{lacroix2020tkbc}: loại bỏ các quan hệ
có dưới \texttt{BUILD\_MIN\_RELATION\_FACTS} sự kiện và các thực thể xuất hiện trong dưới
\texttt{BUILD\_MIN\_ENTITY\_FACTS} sự kiện, nhằm giữ lại phần đồ thị dày đặc và ổn định về mặt thống
kê. Kết quả là một đồ thị tri thức thời gian với khoảng \textbf{238 nghìn} sự kiện, \textbf{37,5
nghìn} thực thể và \textbf{271} quan hệ, trong đó 100\% nhãn cho cặp chủ thể -- đối tượng đều có
phiên bản tiếng Việt. Mỗi sự kiện được lưu dưới dạng một bộ bốn $(s, r, o, \tau)$ với
$\tau = [t_{b}, t_{e}]$ (cho phép $t_{e}$ rỗng với các sự kiện còn hiệu lực).

\subsection{Phủ tiếng Việt và dịch nhãn còn thiếu}
\label{sec:vicoverage}

Khảo sát trên đồ thị thu được cho thấy độ phủ nhãn tiếng Việt không đồng đều: khoảng \textbf{34\%}
thực thể và \textbf{77\%} quan hệ có sẵn nhãn tiếng Việt trên Wikidata. Để bảo đảm mọi câu hỏi sinh
ra đều thuần Việt, phần nhãn còn thiếu được dịch tự động bằng một mô hình ngôn ngữ lớn (Qwen), rồi
lưu thành các tệp nhãn tiếng Việt riêng (\texttt{qid\_labels.vi.json}, \texttt{pid\_labels.vi.json}).
Nhờ vậy, khâu sinh câu hỏi ở giai đoạn sau không cần thao tác dịch nào nữa.

\subsection{Sinh câu hỏi}
\label{sec:generate}

Mô-đun \texttt{generate.py} sinh câu hỏi từ một tập \textbf{21 quan hệ} được tuyển chọn, tổ chức
thành \textbf{8 nhóm chủ đề} (mở rộng từ 5 quan hệ lõi của CronQuestions), như liệt kê trong
Bảng~\ref{tab:relations}. Mỗi quan hệ được gắn một động từ và một danh từ tiếng Việt để ghép vào
mẫu câu (\textit{template}).

\begin{table}[ht]
\centering
\caption{Tám nhóm quan hệ dùng để sinh câu hỏi trong CronQ-VN.}
\label{tab:relations}
\begin{tabular}{L{3.2cm} L{9.5cm}}
\toprule
\textbf{Nhóm} & \textbf{Quan hệ (mã Wikidata)} \\
\midrule
Lõi CronQ & chức vụ (P39), câu lạc bộ thể thao (P54), giải thưởng (P166), hôn nhân (P26), nơi làm việc (P108) \\
Học vấn/Sự nghiệp & học tại (P69), học vị (P512), nơi công tác (P937) \\
Chính trị & đứng đầu chính phủ (P6), đảng phái (P102), chủ tịch (P488), giữ vị trí (P1308) \\
Giải thưởng & đề cử (P1411), thắng cuộc thi (P1346) \\
Thành viên & thành viên của (P463), nơi cư trú (P551) \\
Thể thao & huấn luyện viên trưởng (P286) \\
Quân sự & quân hàm (P410), quân chủng (P241) \\
Khác & chủ sở hữu (P127), hướng dẫn nghiên cứu sinh (P185) \\
\bottomrule
\end{tabular}
\end{table}

Câu hỏi được sinh theo \textbf{năm loại} (mỗi loại chiếm 20\% theo cấu hình mặc định), tương ứng
năm dạng suy luận thời gian, minh họa trong Bảng~\ref{tab:qtype-examples}:

\begin{table}[ht]
\centering
\caption{Năm loại câu hỏi trong CronQ-VN kèm ví dụ.}
\label{tab:qtype-examples}
\begin{tabular}{L{3.0cm} L{9.7cm}}
\toprule
\textbf{Loại} & \textbf{Ví dụ} \\
\midrule
\texttt{simple\_entity} & ``Ai giữ chức Tổng thống Hoa Kỳ vào năm 2010?'' (đáp án: thực thể) \\
\texttt{simple\_time} & ``Barack Obama bắt đầu giữ chức Tổng thống Hoa Kỳ năm nào?'' (đáp án: thời gian) \\
\texttt{before\_after} & ``Ai là Tổng thống Hoa Kỳ ngay trước Barack Obama?'' (đáp án: thực thể) \\
\texttt{first\_last} & ``Câu lạc bộ đầu tiên trong sự nghiệp của Lionel Messi là gì?'' (đáp án: thực thể) \\
\texttt{time\_join} & ``Ai cùng chơi cho FC Barcelona khi Lionel Messi đang ở đó?'' (đáp án: thực thể) \\
\bottomrule
\end{tabular}
\end{table}

Với mỗi câu hỏi, hệ thống ghép mẫu câu rồi \emph{truy vấn trực tiếp trên đồ thị} để thu được tập
đáp án chuẩn (\textit{ground truth}). Chẳng hạn, để sinh câu hỏi \texttt{simple\_entity}, hệ thống
chọn một quan hệ và một năm, lấy tập các sự kiện đang có hiệu lực tại năm đó (\texttt{active\_at}),
và tập chủ thể của chúng chính là đáp án. Cách sinh ``từ đồ thị ra câu hỏi'' này bảo đảm đáp án luôn
nhất quán với tri thức trong đồ thị.

Để bảo đảm chất lượng, một loạt bộ lọc được áp dụng ở khâu sinh:

\begin{itemize}
  \item \textbf{Lọc nhãn xấu:} bỏ nhãn quá ngắn (dưới \texttt{MIN\_LABEL\_LEN} ký tự) hoặc toàn ký
  tự số; bỏ các nhãn đối tượng mang tính lớp chung gây câu hỏi vô nghĩa (ví dụ ``chính trị gia độc
  lập'' cho quan hệ đảng phái).
  \item \textbf{Không tự tham chiếu:} loại sự kiện có chủ thể trùng đối tượng.
  \item \textbf{Không rò rỉ đáp án:} loại câu hỏi mà đáp án xuất hiện ngay trong phần đề câu hỏi.
  \item \textbf{Giới hạn số đáp án:} bỏ câu hỏi có quá \texttt{MAX\_ANSWERS} đáp án (tránh các câu
  ``bùng nổ'' như danh sách nghị sĩ).
  \item \textbf{Ràng buộc khoảng cách thời gian:} câu hỏi \texttt{before\_after} yêu cầu hai sự
  kiện cách nhau tối thiểu \texttt{MIN\_BEFORE\_AFTER\_GAP} năm; câu hỏi \texttt{time\_join} yêu
  cầu hai sự kiện chồng lấn tối thiểu \texttt{TIME\_JOIN\_MIN\_OVERLAP} năm.
  \item \textbf{Khử trùng lặp:} loại các câu hỏi trùng nội dung.
\end{itemize}

\subsection{Lược đồ dữ liệu}
\label{sec:schema}

Mỗi câu hỏi được lưu dưới dạng một bản ghi JSON. Quan trọng nhất, trường \texttt{entities} lưu sẵn
danh sách mã \texttt{QID} của các thực thể mỏ neo, giúp pha suy luận khởi tạo ngay mà không cần
liên kết thực thể lại; còn trường \texttt{template} lưu mẫu câu đã trừu tượng hóa, được dùng cho
việc phân cụm theo cấu trúc (Mục~\ref{sec:knowledge-agnostic}).

\begin{lstlisting}[language=Python,caption={Một bản ghi câu hỏi trong CronQ-VN.},label=lst:schema]
{
  "quid": 1,
  "question": "Ai giữ chức Tổng thống Hoa Kỳ vào năm 2010?",
  "answers": ["Barack Obama"],
  "answer_type": "entity",        # entity | time
  "qtype": "simple_entity",
  "time_level": "year",
  "relation": "P39",
  "year": 2010,
  "entities": ["Q30"],            # QID mỏ neo (seed)
  "template": "Ai {relation} {entity} vào {time}?",
  "qlabel": "Single"              # Single | Multiple
}
\end{lstlisting}

\section{Pha tương tác dựa trên tri thức}
\label{sec:knowledge-based}

Pha này hiện thực hóa toàn bộ tương tác với đồ thị tri thức, theo đúng tinh thần §3.3 của~
\cite{chen2024ari}: biểu diễn bài toán suy luận phức tạp thành chuỗi suy luận nhiều bước, ở mỗi
bước sinh ra một tập hành động ứng viên khả thi để mô hình ngôn ngữ lớn lựa chọn.

\subsection{Đồ thị con và chỉ mục đồ thị}
\label{sec:onehop}

Đồ thị tri thức thời gian được nạp toàn bộ vào bộ nhớ và lập chỉ mục theo chủ thể, đối tượng, quan
hệ và năm để truy vấn nhanh. Xuất phát từ thực thể chủ thể $e_h$ của câu hỏi, hệ thống xét đồ thị
con một bước (\textit{1-hop subgraph}) gồm mọi sự kiện có $e_h$ tham gia:

\begin{equation}
\label{eq:subgraph}
G_{e_h} = \{(e, r) \mid e \in N_{e_h},\ r \in R_{e_h}\},
\end{equation}

trong đó $N_{e_h}$ là tập các nút kề $e_h$ và $R_{e_h}$ là tập cạnh tương ứng. Trong cài đặt, đồ
thị con này chính là hợp của các sự kiện mà $e_h$ đóng vai trò chủ thể hoặc đối tượng.

\subsection{Chín thao tác nguyên tử}
\label{sec:actions}

Thay vì các hàm bậc cao được định nghĩa thủ công, hệ thống dùng một tập \textbf{chín thao tác
nguyên tử} (\textit{atomic operations}) có thể kết hợp linh hoạt để tạo nên các truy vấn phức tạp,
tương ứng Bảng~6 của~\cite{chen2024ari}. Bảng~\ref{tab:actions} liệt kê chín thao tác này.

\begin{table}[ht]
\centering
\caption{Chín thao tác nguyên tử trong hệ thống.}
\label{tab:actions}
\begin{tabular}{L{6.0cm} L{6.7cm}}
\toprule
\textbf{Thao tác} & \textbf{Ý nghĩa} \\
\midrule
\texttt{get\_tail\_entity(head, rel, time)} & Lấy đối tượng từ chủ thể, quan hệ, thời gian \\
\texttt{get\_head\_entity(tail, rel, time)} & Lấy chủ thể từ đối tượng, quan hệ, thời gian \\
\texttt{get\_time(head, rel, tail)} & Lấy thời gian của một sự kiện cụ thể \\
\texttt{get\_between(ents, t1, t2)} & Lọc thực thể có thời gian trong khoảng $[t_1, t_2]$ \\
\texttt{get\_before(ents, t)} & Lọc thực thể có thời gian trước $t$ \\
\texttt{get\_after(ents, t)} & Lọc thực thể có thời gian sau $t$ \\
\texttt{get\_first(ents)} & Lấy thực thể có thời gian sớm nhất \\
\texttt{get\_last(ents)} & Lấy thực thể có thời gian muộn nhất \\
\texttt{answer(value)} & Trả về đáp án cuối cùng \\
\bottomrule
\end{tabular}
\end{table}

Mỗi thao tác trả về một \textit{tập thực thể} -- danh sách các bộ ba (nhãn, mã định danh, năm) --
hoặc một giá trị đáp án. Các thao tác lọc theo thời gian (\texttt{get\_before}, \texttt{get\_after},
\texttt{get\_between}, \texttt{get\_first}, \texttt{get\_last}) nhận đầu vào là kết quả của bước
trước, cho phép \emph{nối chuỗi} các thao tác để tạo nên suy luận nhiều bước. Tính tách rời thấp
giữa hệ thống và tập thao tác giúp dễ dàng mở rộng sang dữ liệu hay miền mới: chỉ cần bổ sung thao
tác nguyên tử tương ứng.

\subsection{Liệt kê và lọc hành động ứng viên}
\label{sec:enumerate}

Ở mỗi bước, hệ thống \emph{liệt kê} mọi hành động khả thi. Từ đồ thị con của các thực thể mỏ neo,
tập hành động khởi tạo được sinh ra bằng cách duyệt và thay thế các quan hệ, thực thể hiện diện:

\begin{equation}
\label{eq:enum}
P_0 = \{\mathrm{Enum}(\textit{action}, e, r) \mid e, r \in G_{e_h}\}.
\end{equation}

Quan trọng, chỉ những hành động \emph{thực thi cho kết quả khác rỗng} mới được giữ lại. Tuy nhiên,
do sự kiện thời gian rất nhiều, ngay cả đồ thị con một bước cũng có thể sinh ra một lượng lớn hành
động, gây nhiễu cho quyết định của mô hình ngôn ngữ lớn. Vì vậy, hệ thống áp dụng một bước
\emph{lọc}: trước hết loại các hành động trùng lặp và các hành động đã được chọn ở bước trước (tránh
lặp vô hạn); sau đó, nếu số ứng viên vẫn vượt ngưỡng \texttt{TOP\_K\_ACTIONS}, giữ lại $K$ hành động
có \emph{độ tương đồng ngữ nghĩa cao nhất} với câu hỏi:

\begin{equation}
\label{eq:filter}
P'_0 = \{a \mid \mathrm{exec}(a) \neq \emptyset \ \wedge\ a \in \mathrm{Top\text{-}K}(P_0, q)\}.
\end{equation}

Độ tương đồng được tính bằng cosine giữa biểu diễn nhúng của \emph{chuỗi mô tả hành động} và của
\emph{câu hỏi}. Như sẽ phân tích ở Mục~\ref{sec:crossencoder}, chính cách lọc ``so hành động với
câu hỏi gốc'' này bộc lộ hạn chế khi suy luận nhiều bước, và là động lực cho đóng góp thứ ba của
khóa luận.

\section{Pha suy luận trừu tượng}
\label{sec:knowledge-agnostic}

Pha suy luận trừu tượng cho phép mô hình ngôn ngữ lớn chắt lọc và áp dụng các phương pháp luận trừu
tượng từ kinh nghiệm lịch sử, theo §3.4 của~\cite{chen2024ari}.

\subsection{Lưu trữ lịch sử suy luận}
\label{sec:history}

Trong quá trình giải mỗi câu hỏi, hệ thống ghi lại trạng thái tại mỗi bước $t_i$, gồm câu hỏi $q$,
tập hành động ứng viên $P_{t_i}$ và quyết định $a_{t_i}$ của mô hình. Tập các trạng thái này tạo
thành lịch sử suy luận của câu hỏi:

\begin{equation}
\label{eq:history}
H_q = \{(q, t_i, P_{t_i}, a_{t_i}) \mid i \in T\},
\end{equation}

với $T$ là tập các bước. Mỗi bản ghi lưu đầy đủ phản hồi của mô hình (cả hành động lẫn lý do) cùng
kết quả truy xuất, phục vụ cho việc chắt lọc phương pháp luận ở giai đoạn sau.

\subsection{Phân cụm theo cấu trúc và chắt lọc phương pháp luận}
\label{sec:induce}

Suy luận thời gian tuy nhiều bước nhưng các \emph{kiểu} suy luận lại khá ổn định: những câu hỏi
tương tự đòi hỏi các bước suy luận tương tự. Do đó, sau khi tích lũy một tập lịch sử, hệ thống dùng
thuật toán K-means (Mục~\ref{sec:kmeans}) để phân các trace thành các cụm.

Một cải tiến quan trọng so với cách làm gốc là \textbf{phân cụm theo cấu trúc thay vì theo chủ đề}.
Trước khi nhúng để phân cụm, mỗi câu hỏi được \emph{trừu tượng hóa}: thay nhãn thực thể, nhãn quan
hệ và mốc thời gian cụ thể bằng các ký hiệu giữ chỗ (ví dụ ``Ai \{relation\} \{entity\} vào
\{time\}?''). Việc thay cả nhãn quan hệ là cần thiết cho tiếng Việt, vì nhãn quan hệ thường trùng
với danh từ trong câu, dễ khiến cụm bị gom theo chủ đề thay vì theo dạng suy luận. Trong cài đặt,
mẫu câu trừu tượng này có sẵn ở trường \texttt{template} của mỗi câu hỏi.

Với mỗi cụm, hệ thống chắt lọc một \emph{phương pháp luận trừu tượng} $M_C$ bằng cách cho mô hình
ngôn ngữ lớn quan sát \emph{cả ví dụ đúng lẫn ví dụ sai} trong cụm (tối đa ba ví dụ mỗi loại) rồi
khái quát thành chỉ dẫn từng bước, độc lập với tri thức cụ thể. Việc đưa cả ví dụ sai giúp mô hình
học cách tránh các lỗi đã gặp, phù hợp với quan sát của~\cite{chen2024ari}. Mỗi cụm cuối cùng lưu
lại tâm cụm, kích thước, phân bố loại câu hỏi và phương pháp luận tương ứng.

\subsection{Suy luận với hướng dẫn trừu tượng}
\label{sec:inference-loop}

Khi gặp câu hỏi mới, hệ thống trừu tượng hóa câu hỏi giống như lúc xây bộ nhớ, rồi chọn cụm có tâm
gần nhất theo khoảng cách Euclid:

\begin{equation}
\label{eq:select}
C^* = \arg\max_{C_i} S(C_i, q),
\end{equation}

và lấy phương pháp luận $M_{C^*}$ của cụm đó. Tại mỗi bước, mô hình ngôn ngữ lớn nhận phương pháp
luận, câu hỏi, lịch sử các bước đã thực hiện và tập hành động ứng viên đã lọc $P'_t$, rồi chọn ra
hành động kế tiếp:

\begin{equation}
\label{eq:decide}
a^*_i = \mathrm{LLM}(M_{C^*},\ q,\ P'_i).
\end{equation}

Hành động được thực thi trên đồ thị, cập nhật trạng thái cho bước sau. Vòng lặp kết thúc khi mô hình
chọn hành động \texttt{answer(\dots)} hoặc khi đạt số bước tối đa
(\texttt{MAX\_REASONING\_STEPS}). Thuật toán~\ref{alg:ari} tóm tắt toàn bộ vòng lặp.

\begin{algorithm}[ht]
\caption{Suy luận trừu tượng quy nạp (vòng lặp suy luận)}
\label{alg:ari}
\begin{algorithmic}[1]
\Require Đồ thị $\mathcal{K}$, câu hỏi $q$, bộ nhớ phương pháp luận $\mathcal{M}$
\Ensure Đáp án cho $q$
\State $M_{C^*} \gets$ chọn phương pháp luận cho $q$ theo Công thức~\eqref{eq:select}
\State Khởi tạo tập thực thể mỏ neo (từ trường \texttt{entities} hoặc liên kết thực thể)
\State $prev \gets \emptyset$
\For{$t = 0$ \textbf{đến} \texttt{MAX\_REASONING\_STEPS}$-1$}
  \State Liệt kê hành động ứng viên $P_t$ từ mỏ neo và $prev$ (Công thức~\eqref{eq:enum})
  \State Lọc còn $P'_t$: bỏ trùng, bỏ đã dùng, giữ Top-K (Công thức~\eqref{eq:filter})
  \State $a^*_t \gets \mathrm{LLM}(M_{C^*}, q, \text{lịch sử}, P'_t)$ \Comment{Công thức~\eqref{eq:decide}}
  \If{$a^*_t$ là \texttt{answer}}
    \State \Return giá trị đáp án
  \EndIf
  \State $prev \gets \mathrm{exec}(a^*_t)$; cập nhật lịch sử
\EndFor
\State \Return ứng viên tốt nhất còn lại trong $prev$
\end{algorithmic}
\end{algorithm}

\section{Thích nghi cho tiếng Việt và vận hành cục bộ}
\label{sec:vietnamese-adapt}

Đóng góp thứ hai của khóa luận là đưa toàn bộ khung ARI sang tiếng Việt và vận hành được trên hạ
tầng cục bộ.

\subsection{Mô hình ngôn ngữ lớn cục bộ}
\label{sec:local-llm}

Khác với bản gốc dùng API thương mại (\texttt{gpt-3.5-turbo}), hệ thống vận hành pha suy luận hoàn
toàn cục bộ thông qua Ollama~\cite{ollama}, với mô hình mặc định Qwen2.5~\cite{qwen2025} (cấu hình
\texttt{qwen3:14b}) và mô hình nhúng \texttt{nomic-embed-text}. Toàn bộ câu lệnh hệ thống và mẫu câu
lệnh được viết bằng tiếng Việt, gồm câu lệnh chọn hành động và câu lệnh chắt lọc phương pháp luận.
Nhiệt độ giải mã được đặt $T = 0$ để bảo đảm khả năng lặp lại (Mục~\ref{sec:llm}).

Khung được thiết kế \emph{trung lập với mô hình}: pha chắt lọc phương pháp luận có thể dùng một mô
hình mạnh hơn để bám sát bản gốc (ví dụ GPT-4o), trong khi pha suy luận vẫn chạy bằng mô hình cục
bộ. Điểm cốt lõi của đóng góp là \emph{quá trình suy luận} -- vốn được gọi lặp nhiều lần và nhạy
cảm về chi phí -- có thể vận hành hoàn toàn cục bộ, không phụ thuộc dịch vụ bên ngoài.

\subsection{Liên kết thực thể cho câu hỏi mở}
\label{sec:entity-linking}

Với các câu hỏi trong bộ CronQ-VN, thực thể mỏ neo đã có sẵn ở trường \texttt{entities}. Tuy nhiên,
để hệ thống trả lời được \emph{câu hỏi mở} do người dùng nhập, cần một bước \emph{liên kết thực thể}
ánh xạ cụm từ trong câu hỏi về đúng mã \texttt{QID} trong đồ thị. Quy trình được thiết kế gồm bốn
bước (Hình~\ref{fig:entity-linking}): (1) nhận câu hỏi; (2) dùng mô hình NER tiếng Việt
(\texttt{NlpHUST/ner-vietnamese-electra-base}) trích xuất các cụm chỉ thực thể; (3) truy vấn các cụm
này vào một cơ sở dữ liệu vector chứa biểu diễn nhúng của nhãn thực thể, tìm theo độ tương đồng
cosine; (4) khớp về các mã \texttt{QID} tương ứng.

\begin{figure}[ht]
  \centering
  \fbox{\parbox{0.93\textwidth}{\centering\vspace{2mm}\footnotesize
  Câu hỏi $\rightarrow$ \textbf{NER} (trích cụm thực thể) $\rightarrow$ \textbf{VectorDB}
  (cosine) $\rightarrow$ khớp \texttt{QID}\\[1mm]
  ``Ai là tổng thống đầu tiên của Hoa Kỳ?'' $\Rightarrow$ \{tổng thống, Hoa Kỳ\} $\Rightarrow$
  \{\texttt{Q30}, \dots\}
  \vspace{2mm}}}
  \caption{Liên kết thực thể cho câu hỏi mở: NER tiếng Việt trích cụm thực thể, sau đó khớp về mã
  định danh qua tìm kiếm tương đồng trong cơ sở dữ liệu vector. \todo{vẽ lại bằng TikZ}}
  \label{fig:entity-linking}
\end{figure}

Trong cài đặt hiện tại, bên cạnh quy trình NER nói trên, hệ thống còn cung cấp một bộ liên kết
\emph{nhẹ} dựa trên so khớp chuỗi con: chuẩn hóa câu hỏi, giới hạn không gian tìm kiếm theo ngữ cảnh
quan hệ của câu hỏi, rồi khớp nhãn thực thể xuất hiện trong câu (ưu tiên nhãn dài nhất) hoặc khi mọi
từ của nhãn đều có mặt trong câu. \todo{xác nhận với Danh: phiên bản dùng để chạy thực nghiệm cuối
cùng là NER + VectorDB hay bộ so khớp chuỗi con, để mô tả cho khớp.}

\section{Cải tiến: bộ lọc hành động bằng cross-encoder}
\label{sec:crossencoder}

Đóng góp thứ ba của khóa luận là một bộ lọc hành động \emph{nhận biết ngữ cảnh} nhằm khắc phục hạn
chế của bộ lọc dựa trên độ tương đồng cosine.

\subsection{Hạn chế của bộ lọc cosine}
\label{sec:cosine-limit}

Như đã nêu ở Mục~\ref{sec:enumerate}, bộ lọc Top-K hiện tại (Công thức~\eqref{eq:filter}) xếp hạng
mỗi hành động dựa trên độ tương đồng cosine giữa \emph{chuỗi mô tả hành động} và \emph{câu hỏi gốc}.
Cách làm này có một khiếm khuyết then chốt: độ tương đồng \emph{không phụ thuộc vào bước suy luận
hiện tại}. Dù hệ thống đang ở bước thứ nhất hay bước thứ tư, thứ hạng của các hành động vẫn giống hệt
nhau, vì đều được so với cùng một câu hỏi gốc mà không xét tới lịch sử các hành động đã thực hiện.

Hệ quả là hành động \emph{đúng cho bước hiện tại} có thể bị xếp hạng thấp hơn một hành động đúng cho
một bước khác, và do đó bị loại nhầm khi cắt Top-K. Đối với suy luận nhiều bước -- vốn là trọng tâm
của bài toán -- đây là một nguồn lỗi đáng kể.

\subsection{Bộ lọc nhận biết ngữ cảnh dựa trên cross-encoder}
\label{sec:ce-filter}

Khóa luận đề xuất thay bước lọc Top-K thuần cosine bằng một quy trình \emph{hai tầng}: dùng cosine
để chọn \emph{rộng} một tập ứng viên (nhằm không bỏ sót hành động đúng), sau đó dùng một
\textbf{cross-encoder} để \emph{xếp hạng lại} và chọn Top-K cuối cùng. Khác với cosine, cross-encoder
mã hóa \emph{đồng thời} hành động và \emph{lịch sử suy luận}, nhờ đó nắm bắt được ngữ cảnh của bước
hiện tại (Mục~\ref{sec:crossencoder-bg}).

Cụ thể, với mỗi hành động ứng viên $a$, bộ chấm điểm nhận đầu vào là chuỗi ghép
$[\mathtt{CLS}]\,a\,[\mathtt{SEP}]\,h\,[\mathtt{SEP}]$, trong đó $h$ là \emph{lịch sử} gồm câu hỏi
và chuỗi các hành động đúng đã thực hiện. Chuỗi này đi qua BERT rồi một lớp tuyến tính kèm sigmoid để
cho ra điểm phù hợp:

\begin{equation}
\label{eq:ce-score}
\mathrm{score}(a, h) = \sigma\!\left(\mathbf{w}^{\top}\,
\mathrm{BERT}_{\mathtt{CLS}}\big([\mathtt{CLS}]\,a\,[\mathtt{SEP}]\,h\,[\mathtt{SEP}]\big) + b\right).
\end{equation}

Nhờ gộp chung biểu diễn của hành động và lịch sử, bộ chấm điểm học được \emph{liên kết giữa hành
động và ngữ cảnh}, qua đó xếp hạng đúng những hành động phù hợp với bước hiện tại -- điều mà cosine
không làm được. Hình~\ref{fig:crossencoder} minh họa kiến trúc.

\begin{figure}[ht]
  \centering
  \fbox{\parbox{0.93\textwidth}{\centering\vspace{2mm}\footnotesize
  Cosine (chọn rộng) $\Rightarrow$ tập ứng viên $\Rightarrow$
  $[\mathtt{CLS}]\,a\,[\mathtt{SEP}]\,h\,[\mathtt{SEP}] \rightarrow \mathrm{BERT} \rightarrow
  \mathrm{Linear} \rightarrow \sigma \rightarrow \mathrm{score} \rightarrow$ Top-K cuối\\[1mm]
  với $h$ = câu hỏi + các hành động đúng đã thực hiện
  \vspace{2mm}}}
  \caption{Bộ lọc hành động hai tầng đề xuất: cosine chọn rộng, cross-encoder xếp hạng lại theo ngữ
  cảnh lịch sử. \todo{vẽ lại bằng TikZ}}
  \label{fig:crossencoder}
\end{figure}

\subsection{Huấn luyện bộ chấm điểm}
\label{sec:ce-train}

Bộ chấm điểm được huấn luyện như một bộ phân loại nhị phân với hàm mất mát entropy chéo nhị phân
(Công thức~\eqref{eq:bce}). Dữ liệu huấn luyện được sinh từ các trace suy luận lịch sử: với mỗi
bước, hành động \emph{thực sự dẫn tới đáp án đúng} được gán nhãn dương, còn các hành động ứng viên
khác (đặc biệt là các hành động đúng cho bước \emph{khác} nhưng sai cho bước hiện tại) được gán nhãn
âm. Cách lấy mẫu âm ``khó'' này buộc mô hình học ranh giới phân biệt theo ngữ cảnh.
\todo{bổ sung cấu hình huấn luyện chi tiết (mô hình BERT nền, số epoch, kích thước batch, tỉ lệ mẫu
âm) sau khi triển khai và chạy.}

\section{Tổng kết chương}

Chương này đã trình bày đầy đủ hệ thống TKGQA tiếng Việt của khóa luận: kiến trúc hai pha kế thừa
ARI, quy trình xây dựng bộ dữ liệu CronQ-VN, pha tương tác tri thức với chín thao tác nguyên tử, pha
suy luận trừu tượng dựa trên phân cụm và chắt lọc phương pháp luận, các thích nghi cho tiếng Việt và
vận hành cục bộ, cùng bộ lọc hành động dựa trên cross-encoder. Chương tiếp theo trình bày thiết lập
thực nghiệm và đánh giá hệ thống theo ba câu hỏi nghiên cứu đã đặt ra.
